\documentclass[10pt,twocolumn,letterpaper]{article}

\usepackage{statcourse}
\usepackage{times}
\usepackage{epsfig}
\usepackage{graphicx}
\usepackage{amsmath}
\usepackage{amssymb}

% Include other packages here, before hyperref.

% If you comment hyperref and then uncomment it, you should delete
% egpaper.aux before re-running latex.  (Or just hit 'q' on the first latex
% run, let it finish, and you should be clear).
\usepackage[breaklinks=true,bookmarks=false]{hyperref}


\statcoursefinalcopy


\setcounter{page}{1}
\begin{document}


%%%%%%%%%%%%%%%%%%%%%%%%%%%%%%%%%%%%%%%%%%%%%%%%%%%%%%%%%%%%%%%
% DO NOT EDIT ANYTHING ABOVE THIS LINE
% EXCEPT IF YOU LIKE TO USE ADDITIONAL PACKAGES
%%%%%%%%%%%%%%%%%%%%%%%%%%%%%%%%%%%%%%%%%%%%%%%%%%%%%%%%%%%%%%%



%%%%%%%%% TITLE
\title{\LaTeX\ STAT479 Project Proposal}

\author{Tong Li\\
{\tt\small tli287@wisc.edu}
\and
Runxin Gao\\
{\tt\small secondauthor@wisc.edu}
\and
Jianian Jin\\
{\tt\small thirdauthor@wisc.edu}
}

\maketitle
%\thispagestyle{empty}



% MAIN ARTICLE GOES BELOW
%%%%%%%%%%%%%%%%%%%%%%%%%%%%%%%%%%%%%%%%%%%%%%%%%%%%%%%%%%%%%%%



%%%%%%%%% BODY TEXT



\section{Introduction}


In this section, describe what you are planning to do. Also, briefly describe related work.


When discussing related work, do not forget to include appropriate references.  This is an example of a citation \cite{mirjalili2018gender}. To format the citations properly, put the
corresponding references into the bibliography.bib file. You can obtain
BibTeX-formatted references for the "bib" file from Google Scholar 
(\url{https://scholar.google.com}), for example, by clicking on the 
double-quote character under a citation and then selecting \mbox{"BibTeX"} as
shown in Figure \ref{fig:google-scholar-1col} and 
Figure \ref{fig:google-scholar-2col}.


\begin{figure}[t]
\begin{center}
   \includegraphics[width=0.8\linewidth]{figures/quickdraw_image.jpg}
\end{center}
   \caption{Example of quick draw imgaes}
\label{fig:quick draw image}
\end{figure}


\begin{figure*}
\begin{center}
   \includegraphics[width=0.8\linewidth]{figures/google-scholar.pdf}
\end{center}
   \caption{Example illustrating how to get BibTeX references from
   Google Scholar as a 2-column figure.}
\label{fig:google-scholar-2col}
\end{figure*}


\section{Motivation}

Describe why your project is interesting. E.g., you can describe why your project could have a broader societal impact. Or, you may describe the motivation from a personal learning perspective.

\section{Evaluation}

What would the successful outcome of your project look like? In other words, under which circumstances would you consider your project to be “successful?”

How do you measure success, specific to this project, from a technical standpoint?

\section{Resources}

Google has made its Quick! Draw dataset public with several versions, such as the original raw data, the simplified version of the raw data, as well as numpy bitmap data, built upon the simplified data, that converts the data into bitmaps with $28x28$ grayscale which has 784 dimensions. Because the aim is to classify drawings from Quick! Draw, we’ve decided to use to numpy bitmap data.


This dataset consists of 345 categories and over fifty million drawings in total, which would make the computation time really long. Thus, we are planning to draw a random sample from each category to make computation faster. Then we would divide the sample into training, validation and test set for our model construction and evaluation.


\section{Contributions}

You are expected to share the workload evenly, and every group member is expected to participate in both the experiments and writing. (As a group, you only need to submit one proposal and one report, though. So you need to work together and coordinate your efforts.)

Clearly indicate what computational and writing task each member of your group will be participating in.


{\small
\bibliographystyle{ieee}
\bibliography{bibliography.bib}
}

\end{document}
